\section{Discussion and Threats to Validity}

\subsection{What the Generation Check Protects}

The controlled interleaving distinguishes two responses. The delayed GET holds a PostgreSQL snapshot that was valid when its query ran. The later GET starts after the final order has acknowledged. Mode C can return the older value to the later GET because the delayed request repopulates a key that the writer has already deleted. In D, the generation changes while the delayed request is paused. Redis rejects the refill even when the payload key is absent. The later GET can obtain a newer snapshot. This is a property of the exercised schedule under a successful Redis invalidation. The original overlapping GET may still carry its older snapshot.

The earlier sequential probe primed a cache key and read it again soon after the order, inside the ten-second TTL. That sequence illustrates expiration semantics and conditions on cache priming. Even if 15 probes were independent Bernoulli trials, zero stale responses would give a two-sided 95\% exact upper bound of \(1-0.025^{1/15}=21.8\%\) for the stale probability \cite{clopper1934}. The controlled replay and the concurrent stress test answer different questions. The replay forces one interleaving. The stress test measures naturally scheduled overlaps for one item and one client process. Reads share state within each trial, and all of the test conditions differ from the mix of interactions at a real restaurant.

D uses extra Redis work per read miss and invalidation. A hot item's generation survives payload expiry and consumes Redis memory. If Redis loses generations, times out during invalidation, or serves an operation across a restart, the tested fence assumptions can fail. PostgreSQL and Redis commit independently, so a reader can also observe a stale hit before an invalidation has applied. A durable transition queue and a cache reset or epoch policy would be useful extensions for an environment that requires recovery across service failures.

\subsection{Operational Interpretation}

The depletion trial dataset puts the user-facing effect in observable terms. With one portion displayed and zero portions in PostgreSQL, a browser following the display submits a checkout that the server must reject. These outcomes show unnecessary attempts and do not identify the rate of lost sales. Eight actual clients share one depletion event, real customers may browse differently, and a real menu may be restocked. A longer study across multiple items and stock replenishments would be needed to estimate routine operational exposure.

The performance comparison combines fixed offered rates, per-endpoint timing, and in-window delivery. For a loaded system, a good P95 on successful GETs can coexist with dropped starts and slower POSTs. Counts of item SELECT calls explain which policies reduce application database traffic. CPU, lock wait, and per-command latency would require separate instrumentation. The driver and services shared one host, and the fixed 128-VU allocation can leave requests unscheduled when responses take longer. A separate load-generator host, longer measurement windows, and higher-rate repeat blocks would help isolate the backend's saturation point.

The low-rate TTL sweep isolates a narrow range of settings in two policies. Its 12-second windows capture short bursts and little more than one expiration period at the longest TTL. A longer observation could reveal how repeated expiry and fill cycles change both latency and cache hit rates. The study uses actual trial traffic, one Express process, one PostgreSQL primary, and one Redis server. Results apply to those conditions and the included workload traces.
